\subsection{کالبدشکافی ماژول \lr{factorised\_matrix\_state}}

\subsubsection*{۱. فلسفه طراحی، پارادایم معماری و تمثیل ملموس}
معماری \lr{factorised\_matrix\_state} پاسخی به این چالش بنیادین است: \textbf{چگونه می‌توان ظرفیت زبانی را بالا نگه داشت و هزینه‌های محاسباتی و بن‌بست حافظه را مهار کرد؟}

\textbf{تمثیل ملموس و مثال کاربردی:}
برای درک معماری \lr{factorised\_matrix\_state}، ترکیب رنگ‌های پایه‌ای روی پالت یک نقاش مینیاتوریست چیره‌دست را تجسم کنید. نقاش برای خلق هزاران رنگ و سایه‌روشن مختلف در تابلوی خود، نیازی ندارد هزار قوطی رنگ مجزا همراه داشته باشد؛ بلکه تنها با ترکیب درصدهای معینی از سه رنگ اصلی، طیف بی‌نهایتی از رنگ‌ها را روی بوم می‌آفریند. این ماژول حالت پنهان بزرگ ماتریسی را به حاصل‌ضرب دو بردار با رتبه پایین $U V^T$ تجزیه کرده و بدون نیاز به حافظه ماتریسی عظیم، پویایی حالت‌های پنهان مرتبه بالا را شبیه‌سازی می‌کند.

\subsubsection*{۲. فرمول‌بندی و استنتاج عمیق ریاضی \lr{(Mathematical Derivation)}}
این مکانیزم با نگاشت خطی تانسورهای پرس‌وجو و کلید از طریق توابع فیچر غیرخطی $\phi(\cdot)$ فرمول‌بندی می‌شود:
\begin{equation}
\operatorname{Attn}(\mathbf{Q}, \mathbf{K}, \mathbf{V}) = \frac{\phi(\mathbf{Q}) \left(\phi(\mathbf{K})^T \mathbf{V}\right)}{\phi(\mathbf{Q}) \sum_{j=1}^N \phi(\mathbf{K})_j^T}
\end{equation}
با تکیه بر ویژگی شرکت‌پذیری ضرب، ماتریس کلید و مقدار به شکل تجمعی $\mathbf{S} \in \mathbb{R}^{d \times d}$ ادغام شده و هزینه به اندازه دنباله مستقل می‌شود.

\subsubsection*{۳. کالبدشکافی الگوریتمی و جریان داده \lr{(Algorithmic Tensor Trace)}}
در پیاده‌سازی سورس پایتون در فایل ماژول \lr{factorised\_matrix\_state.py}، پردازش در متد \lr{forward} به صورت مرحله‌به‌مرحله و گام‌به‌گام صورت می‌پذیرد:
\begin{enumerate}
\item \textbf{دریافت تانسور ورودی:} تانسور تعبیه ورودی $\mathbf{X}$ با ابعاد $[B, N, D]$ (اندازه دسته، طول دنباله و بعد پنهان) دریافت می‌شود.
\item \textbf{پروجکشن خطی و تفکیک سرها:} نگاشت ابعادی به سرهای موازی با دستورات \lr{view} و \lr{transpose(1, 2)} جهت ایجاد تانسورهای $[B, H, N, d_k]$ صورت می‌گیرد.
\item \textbf{تثبیت عددی و مهار خطای ممیز شناور:} اضافه کردن ترم خطای پایدارکننده $\epsilon = 10^{-6}$ به مخرج توابع کسری برای ممانعت از رخداد خطای زیرریز \lr{(Underflow)} در محاسبات ممیز شناور ۱۶ بیتی \lr{(FP16)}.
\item \textbf{کاهش ابعاد و تجمیع اطلاعات:} اجرای عملگر تخصصی معماری با حذف درایه‌های زائد و استخراج مؤلفه‌های اصلی اطلاعات.
\item \textbf{بازگشت به بعد اصلی و پروجکشن پایانی:} تانسور به شکل $[B, N, D]$ بازآرایی شده و با یک لایه خطی پایانی به لایه پسین تزریق می‌گردد.
\end{enumerate}

\subsubsection*{۴. تحلیل همجوشی سخت‌افزاری و ترافیک حافظه \lr{(Hardware \& Roofline Analysis)}}
از دیدگاه مدل سقف کارایی \lr{(Roofline Model)} و معماری پردازنده‌های گرافیکی مدرن (نظیر کارت گرافیک \lr{NVIDIA RTX 4070}):
\begin{itemize}
\item \textbf{مهار عطش پهنای باند حافظه } \lr{(Memory-Bound Relief)}\rl{:}  در مدل‌های استاندارد، تبادل دائم داده‌های ماتریس توجه با حافظه سراسری کند \lr{DRAM/HBM} باعث معطلی هسته‌های محاسباتی می‌شود. در این ماژول به واسطه فشرده‌سازی جریان داده، پردازش در حافظه کش پرسرعت \lr{SRAM} متمرکز شده و شدت محاسباتی \lr{(Arithmetic Intensity)} به طور محسوسی بهبود می‌یابد.
\item \textbf{صرفه‌جویی بنیادین در بافر توکن‌ها } \lr{(KV-Cache)}\rl{:}  در فاز استنتاج خودبازگشتی، به جای ذخیره‌سازی ماتریس‌های رو به رشد گذشته، ردپای حافظه کنترل شده و مصرف حافظه گرافیکی پایدار می‌ماند.
\end{itemize}

\subsubsection*{۵. ارزیابی تجربی، بنچمارک و توجیه علمی نتایج}
بر اساس آزمون‌های ثبت‌شده در خط لوله ارزیابی پایان‌نامه:
\begin{itemize}
\item \textbf{وضعیت رسمی آزمون کیفی:} \lr{PASS}
\item \textbf{دقت ثبت‌شده روی بنچ‌مارک:} \lr{77.06\%}
\item \textbf{شاخص ارزیابی آماری } \lr{(F1 Score)}\rl{:}  \lr{76.08\%}
\item \textbf{اوج مصرف حافظه گرافیکی } \lr{(Peak VRAM)}\rl{:}  \lr{342.5 MB} (در برابر \lr{۴۲۰۰ MB} مدل مبنا)
\item \textbf{نرخ توان عملیاتی } \lr{(Throughput)}\rl{:}  \lr{37,600} توکن بر ثانیه
\end{itemize}

\textbf{توجیه علمی و علت‌شناسی جذاب نتایج:}
دقت \lr{77.06\%}ی و F1 برابر \lr{76.08\%} همراه با مصرف VRAM معادل \lr{342.49 MB} (کاهش \lr{91.8\%}ی در مقایسه با \lr{4200 MB} مبنا)، نشان‌دهنده اثربخشی فشرده‌سازی لایه‌ای بردار حالت در مدل‌های پیوسته است. ماتریس فاکتورشده حافظه لازم برای عملیات اسکن را در محدوده امن نگه داشته است.

\subsubsection*{۶. دیاگرام معماری شکیل در لاتک \lr{(TikZ Architecture Flowchart)}}
جهت تبیین دقیق جریان داده تانسوری و نحوه تغییر شکل ماتریس‌ها، نمودار معماری این ماژول با استفاده از بسته قدرتمند \lr{TikZ} در محیط لاتک ترسیم شده است:

\begin{center}
\begin{tikzpicture}[scale=0.82, every node/.style={scale=0.82}, 
    node distance=1.1cm,
    auto,
    block/.style={rectangle, draw=blue!70, fill=blue!12, rounded corners=4pt, text centered, minimum height=0.85cm, font=\small},
    kernel/.style={rectangle, draw=orange!80, fill=orange!18, rounded corners=4pt, text centered, minimum height=0.85cm, font=\small},
    outblock/.style={rectangle, draw=teal!80, fill=teal!15, rounded corners=4pt, text centered, minimum height=0.85cm, font=\small}
]
    \node [block] (input) {ورودی $\mathbf{X} \in \mathbb{R}^{B \times N \times D}$};
    \node [kernel, right=1.2cm of input] (k) {نگاشت کرنل $\phi(\mathbf{Q}), \phi(\mathbf{K})$};
    \node [block, right=1.2cm of k] (acc) {تجمع ماتریس حافظه $\mathbf{S}_t$};
    \node [outblock, right=1.2cm of acc] (out) {خروجی خطی};

    \draw [->, thick, draw=blue!80] (input) -- (k);
    \draw [->, thick, draw=orange!80] (k) -- (acc);
    \draw [->, thick, draw=blue!80] (acc) -- (out);
\end{tikzpicture}
\end{center}

\subsubsection*{۷. پیوست تصویری و پرامپت تولید هوش مصنوعی}
برای ترسیم تجسم سه‌بعدی و گرافیک صنعتی این مکانیزم با ابزارهای هوش مصنوعی، پرامپت انگلیسی زیر با دقت تدوین شده است:

\vspace{0.2cm}
\begin{latin}
\footnotesize
\textbf{AI Technical Visualization Prompt:}\\
"A hyper-detailed, photorealistic 3D scientific isometric schematic of the factorised\_matrix\_state attention mechanism. Illustrating incoming token tensors of shape [B, N, D] flowing smoothly into optimized projection blocks and state matrices. Glowing cyan, neon amber, and vibrant violet data pathways highlighting low-rank dimension reduction and on-chip SRAM cache operations, contrasted against a dark metallic blueprint background. Clean technical typography, publication-grade academic rendering, Unreal Engine 5 cinematic lighting, 8k resolution --ar 16:9"
\end{latin}


